
\documentclass[11pt]{article}

\usepackage[margin=1in]{geometry}
\usepackage[T1]{fontenc}
\usepackage{lmodern}

\usepackage{amsmath, amssymb}
\usepackage{amsfonts}

\usepackage{setspace}

\usepackage[most]{tcolorbox}
\usepackage{xcolor}

\tcbset{
  enhanced,
  boxrule=0pt,
  frame hidden,
  breakable
}

\usepackage{fancyhdr}
\pagestyle{fancy}
\fancyhf{}
\setlength{\headheight}{52pt}
\renewcommand{\headrulewidth}{0pt}

\newcommand{\Course}{Math 4610 Spring 2026}
\newcommand{\Instructor}{Homework 5}
\newcommand{\AssignmentType}{Homework 5}
\newcommand{\AssignmentNumber}{}

\newcommand{\Contributors}{
Marks, Stephen
}

\fancyhead[L]{\Course\\ \Instructor}
\fancyhead[R]{\Contributors}

\newcommand{\ProblemDivider}[1]{%
  \vspace{1.0em}
  \noindent\textbf{#1}\par
  \vspace{0.35em}
}

\definecolor{problemBack}{RGB}{235,242,250}
\definecolor{problemFrame}{RGB}{70,110,160}
\definecolor{exerciseGray}{gray}{0.96}
\definecolor{exerciseGrayFrame}{gray}{0.45}

\newtcolorbox{problemBox}{
  colback=exerciseGray,
  colframe=exerciseGrayFrame,
  arc=2mm,
  left=8pt,right=8pt,top=8pt,bottom=8pt
}

\newtcolorbox{solutionBox}{
  colback=white,
  colframe=white,
  arc=0mm,
  left=0pt,right=0pt,top=0pt,bottom=0pt
}

\newcommand{\ProblemAnnounce}{\noindent}
\newcommand{\SolutionAnnounce}{\noindent\textbf{Solution.}\par\medskip}

\newcommand{\Exercise}[1]{\textbf{Exercise~#1.}\;}

\newcommand{\R}{\mathbb{R}}
\newcommand{\C}{\mathbb{C}}
\newcommand{\F}{\mathbb{F}}
\newcommand{\id}{\mathrm{id}}

\DeclareMathOperator{\Range}{range}
\DeclareMathOperator{\Null}{null}
\DeclareMathOperator{\Span}{span}

\begin{document}

\begin{center}
  {\Large \textbf{\AssignmentType~\AssignmentNumber}}
\end{center}
\vspace{0.75em}

\ProblemDivider{1}
\begin{problemBox}
\ProblemAnnounce
\Exercise{5.D.2}
Suppose \(T \in \mathcal{L}(V)\) has a diagonal matrix \(A\) with respect to some basis of \(V\). Prove that if \(\lambda \in \mathbf{F}\), then \(\lambda\) appears on the diagonal of \(A\) precisely \(\dim E(\lambda, T)\) times.
\end{problemBox}

\begin{solutionBox}
\SolutionAnnounce
\doublespacing
Let \((v_1,\dots,v_n)\) be a basis of \(V\) such that \(\mathcal{M}(T,(v_1,\dots,v_n))\) is diagonal with diagonal entries \(\lambda_1,\dots,\lambda_n\). For any \(\mu \in \F\), let \(d_\mu\) be the number of times \(\mu\) appears among \(\lambda_1,\dots,\lambda_n\).

First, we show \(\dim E(\mu,T) \le d_\mu\). Consider \(T - \mu I\). In the same basis, \(\mathcal{M}(T-\mu I)\) is diagonal with entries \(\lambda_i - \mu\). The number of zero entries on the diagonal is exactly \(d_\mu\). Since the null space consists of vectors whose coordinates are zero except possibly those corresponding to zero diagonal entries, we have \(\dim \null(T-\mu I) \le d_\mu\). But \(\null(T-\mu I) = E(\mu,T)\), so \(\dim E(\mu,T) \le d_\mu\).

By 5.55, the eigenspaces of \(T\) form a direct sum, so 
\[
\sum_{\mu \in \F} \dim E(\mu,T) = n.
\]
And \(\sum_{\mu \in \F} d_\mu = n\) as well. Since \(\dim E(\mu,T) \le d_\mu\) for all \(\mu\), equality must hold for every \(\mu\). Thus \(\dim E(\lambda,T) = d_\lambda\) for each \(\lambda \in \F\).
\end{solutionBox}

\ProblemDivider{2}

\begin{problemBox}
\ProblemAnnounce
\Exercise{5.D.3}
Suppose \(V\) is finite-dimensional and \(T \in \mathcal{L}(V)\). Prove that if the operator \(T\) is diagonalizable, then \(V = \operatorname{null} T \oplus \operatorname{range} T\).
\end{problemBox}

\begin{solutionBox}
\SolutionAnnounce
\doublespacing
Since \(T\) is diagonalizable,
\[
V = E(0, T) \oplus E(\lambda_1, T) \oplus \dots \oplus E(\lambda_m, T)
\]

where \(\lambda_i\) are the distinct nonzero eigenvalues. Let

\[
W = E(\lambda_1, T) \oplus \dots \oplus E(\lambda_m, T)
\]

so \(V = \Null T \oplus W\).

Take any \(v = u + w_1 + \dots + w_m \in V\), with \(u \in \Null T\) and \(w_i \in E(\lambda_i, T)\). Then

\[
Tv = \lambda_1 w_1 + \dots + \lambda_m w_m \in W.
\]

Thus \(\Range T \subseteq W\). Conversely, for any \(w_1 + \dots + w_m \in W\),

\[
w_1 + \dots + w_m =
T\left(
\frac{w_1}{\lambda_1} + \dots + \frac{w_m}{\lambda_m}
\right)
\in \Range T.
\]

Thus \(W \subseteq \Range T\), so \(W = \Range T\). Therefore

\[
V = \Null T \oplus \Range T.
\]
\end{solutionBox}

\ProblemDivider{3}

\begin{problemBox}
\ProblemAnnounce
\Exercise{5.D.6}
Suppose \(T \in \mathcal{L}(\mathbf{F}^5)\) and \(\dim E(8, T) = 4\). Prove that \(T - 2I\) or \(T - 6I\) is invertible.
\end{problemBox}

\begin{solutionBox}
\SolutionAnnounce
\doublespacing
We will prove the contrapositive: if neither \(T-2I\) nor \(T-6I\) is invertible, then \(\dim E(8, T) \le 3\).

If \(T-2I\) is not invertible, then \(2\) is an eigenvalue, so \(\dim E(2, T) \ge 1\). Similarly, \(\dim E(6, T) \ge 1\).

By 5.54,

\[
\dim E(8, T) + \dim E(2, T) + \dim E(6, T) \le 5.
\]

Then

\[
\dim E(8, T) \le 5 - 1 - 1 = 3.
\]

Thus \(\dim E(8, T) \neq 4\).
\end{solutionBox}

\ProblemDivider{4}

\begin{problemBox}
\ProblemAnnounce
\Exercise{5.D.7}
Suppose \(T \in \mathcal{L}(V)\) is invertible. Prove that

\[
E(\lambda, T) = E(1/\lambda , T^{-1})
\]

for every \(\lambda \in \mathbf{F}\) with \(\lambda \neq 0\).
\end{problemBox}

\begin{solutionBox}
\SolutionAnnounce
\doublespacing
For any nonzero \(\lambda\) and any \(v \in V\),

\[
\begin{aligned}
v \in E(\lambda, T)
&\iff Tv = \lambda v \\
&\iff T^{-1}(Tv) = T^{-1}(\lambda v) \\
&\iff v = \lambda T^{-1}v \\
&\iff T^{-1}v = \frac{1}{\lambda} v \\
&\iff v \in E(1/\lambda, T^{-1}).
\end{aligned}
\]

Thus the eigenspaces are equal as sets and thus as subspaces.
\end{solutionBox}

\ProblemDivider{5}

\begin{problemBox}
\ProblemAnnounce
\Exercise{5.D.11}
Find \(T \in \mathcal{L}(\mathbf{C}^3)\) such that \(6\) and \(7\) are eigenvalues of \(T\) and such that \(T\) does not have a diagonal matrix with respect to any basis of \(\mathbf{C}^3\).
\end{problemBox}

\begin{solutionBox}
\SolutionAnnounce
\doublespacing
Let \(T\) be the operator on \(\C^3\) whose matrix with respect to the standard basis is

\[
\mathcal{M}(T) =
\begin{pmatrix}
6 & 1 & 0 \\
0 & 6 & 0 \\
0 & 0 & 7
\end{pmatrix}.
\]

Because the matrix is upper-triangular, its eigenvalues are \(6\) and \(7\).

\[
E(6, T) = \Span(e_1), \qquad
E(7, T) = \Span(e_3).
\]

Thus

\[
\dim E(6,T) + \dim E(7,T) = 2 \neq 3.
\]

So by 5.55, \(T\) is not diagonalizable.
\end{solutionBox}

\ProblemDivider{Extra Credit}

\begin{problemBox}
\ProblemAnnounce
\textbf{Extra Credit.} Is the converse of Problem 3 valid? Consider \(T(w,z)=(w+z,z)\) on \(V=\mathbf{F}^2\).
\end{problemBox}

\begin{solutionBox}
\SolutionAnnounce
\doublespacing
The converse is not valid. The operator \(T: \F^2 \to \F^2\) defined by
\[
T(w, z) = (w+z, z)
\]
is a counterexample.

First we check the condition \(V = \Null T \oplus \Range T\).

\begin{itemize}
\item The null space:
\[
\Null T = \{(w,z) : w+z = 0,\; z = 0\} = \{(0,0)\}.
\]

\item The range:
\[
\Range T = \{(w+z, z) : w,z \in \F\}.
\]
Let \(x = w+z\) and \(y = z\). Then every \((x,y) \in \F^2\) can be written in this form, so
\[
\Range T = \F^2.
\]
\end{itemize}

Thus
\[
V = \{0\} \oplus \F^2 = \Null T \oplus \Range T.
\]

However, \(T\) is not diagonalizable. The matrix of \(T\) with respect to the standard basis is

\[
\mathcal{M}(T) =
\begin{pmatrix}
1 & 1 \\
0 & 1
\end{pmatrix}.
\]

Its only eigenvalue is \(1\), and the corresponding eigenspace is

\[
E(1,T) = \{(w,0) : w \in \F\} = \Span(e_1),
\]

which has dimension \(1\). Since \(\dim V = 2\), the sum of the dimensions of the eigenspaces is \(1 \neq 2\). By 5.55, \(T\) is not diagonalizable.

Therefore the condition \(V = \Null T \oplus \Range T\) does not imply that \(T\) is diagonalizable.
\end{solutionBox}

\end{document}
